\section{Результаты}
\label{sec:results}

Сначала оценим ускорение симулятора, затем выделим локальные расхождения
и проверим их влияние на выходы и градиенты компонентов.
Для всех численных сравнений используются пары с рис.~\ref{fig:configs}.
Приведённые диапазоны отражают различия слоёв или компонентов;
статистические интервалы не оценивались.

\subsection{Ускорение исполнения симулятора}

Ускорение определим как отношение показаний симулятора
для последовательной реализации на CPU и реализации на GPU
при одинаковых прямом и обратном проходах.
Компиляция XLA в измерение не входит.
В таблице~\ref{tab:speedup} показаны округлённые наружу диапазоны,
охватывающие результаты cuBLAS и CUDA-ядра для каждой модели.
Границы округлены вниз и вверх до десятых; это не доверительные интервалы.

\begin{table}[htbp]
    \centering
    \caption{Ускорение GPU-реализаций относительно последовательной CPU-реализации.}
    \label{tab:speedup}
    \begin{tabular}{l c}
        \hline
        Модель & Диапазон ускорения, раз \\
        \hline
        Трансформер (2 слоя) & $1.9$--$2.0$ \\
        Полносвязная сеть (6 слоёв) & $3.7$--$3.9$ \\
        \hline
    \end{tabular}
\end{table}

\subsection{Послойный профиль и сравнение путей исполнения}

После проверки времени определим, какие операции дают численные
различия. В таблице~\ref{tab:rope} показан фрагмент послойного сравнения 1
для блока с механизмом внимания и вращательным позиционным кодированием (RoPE).
При изолированном повторе расхождения линейных проекций находятся
на уровне $10^{-3}$, тогда как для первого RMSNorm получено
$7.98\cdot10^{-8}$. Таким образом, локальная ошибка зависит от операции.
При позиционном кодировании и случайном занулении
результаты двух реализаций в этом запуске совпали.

Матрица $QK^T$ содержит оценки сходства запросов и ключей.
Чтобы позиция последовательности не использовала последующие позиции,
соответствующие оценки заменяются значением $-10^9$:
это действие каузальной маски. Нормирование экспонент оценок до единичной суммы в каждой строке
(softmax) даёт матрицу весов $A$.
Произведение $AV$ объединяет значения с этими весами.

\begin{table}[htbp]
    \centering
    \caption{Локальные относительные расхождения $d_{\mathrm{rel}}$
        блока с RoPE для сравнения 1, фрагмент.}
    \label{tab:rope}
    \begin{tabular}{p{6.3cm} c c}
        \hline
        Модуль & \shortstack{Прямой\\проход} & \shortstack{Обратный\\проход} \\
        \hline
        Первая нормализация RMSNorm & $7.98\cdot10^{-8}$ & $1.05\cdot10^{-7}$ \\
        Линейные проекции $Q$, $K$, $V$ & $(2.206\text{--}2.234)\cdot10^{-3}$ & $(2.203\text{--}2.242)\cdot10^{-3}$ \\
        Позиционное кодирование RoPE & $0$ & $0$ \\
        Маскированные оценки сходства & $\approx0$ & $2.339\cdot10^{-3}$ \\
        Взвешенная сумма значений $AV$ & $2.331\cdot10^{-3}$ & $2.337\cdot10^{-3}$ \\
        Выходная линейная проекция & $2.250\cdot10^{-3}$ & $2.222\cdot10^{-3}$ \\
        Проекции полносвязного блока & $(2.222\text{--}2.301)\cdot10^{-3}$ & $(2.225\text{--}2.328)\cdot10^{-3}$ \\
        Случайное зануление & $0$ & $0$ \\
        \hline
    \end{tabular}
\end{table}

В строке с маской значения $-10^9$ увеличивают норму опоры
в знаменателе. По незамаскированным позициям получено
$d_{\mathrm{rel}}\approx2.3\cdot10^{-3}$.
Поэтому метрика около нуля здесь не означает совпадения всех элементов.

Послойный профиль ещё не отделяет замену GEMM от изменения
цепочки компиляции. Для этого на рис.~\ref{fig:hierarchy} сведены
локальные расхождения для всех четырёх пар реализаций.

\begin{figure}[htbp]
    \centering
    \begin{tikzpicture}[font=\normalsize]
        \draw[-{Stealth[length=4pt]}, line width=0.6pt]
            (0,0) -- (11.15,0) node[right] {$d_{\mathrm{rel}}$};
        \foreach \x/\power in {0/-8,1.8/-7,3.6/-6,5.4/-5,7.2/-4,9/-3,10.8/-2} {
            \draw (\x,0.07) -- (\x,-0.07);
            \node[below] at (\x,-0.12) {$10^{\power}$};
        }
        \draw[green!45!black, line width=0.7pt] (1.71764,0) -- (1.71764,0.65);
        \fill[green!45!black] (1.71764,0) circle (2pt);
        \node[above, align=center] at (1.71764,0.65)
            {Пара 2\\$\sim9\cdot10^{-8}$};
        \draw[blue!65!black, line width=0.7pt] (5.54253,0) -- (5.54253,0.65);
        \fill[blue!65!black] (5.54253,0) rectangle ++(0.06,0.06);
        \node[above, align=center] at (5.54253,0.65)
            {Пары 3 и 4\\$\sim1.2\cdot10^{-5}$};
        \draw[red!70!black, line width=0.7pt] (9.61636,0) -- (9.61636,0.65);
        \fill[red!70!black] (9.61636,0) circle (2pt);
        \node[above, align=center] at (9.61636,0.65)
            {Пара 1\\$\sim2.2\cdot10^{-3}$};
    \end{tikzpicture}
    \caption{Характерные локальные расхождения GEMM для пар с рис.~\ref{fig:configs}.
        Шкала логарифмическая; значения для пар 3 и 4 совпадают
        с указанной точностью.}
    \label{fig:hierarchy}
\end{figure}

Различия внутри общей цепочки XLA на рис.~\ref{fig:hierarchy}
примерно в $180$ раз (CUDA) и $2.4\cdot10^4$ раз (cuBLAS) ниже
расхождения с CPU-эталоном без XLA. Таким образом, замена реализации
GEMM внутри симулятора
не объясняет весь масштаб различия с исходной программой.
Нужна также проверка общей цепочки исполнения; для установления
конкретного механизма потребовались бы отдельные измерения
слияния операций, порядка суммирования и внутренних форматов.

\subsection{Знаки и корреляция отклонений}

Чтобы дополнить сравнение по норме, рассмотрим распределение знаков
отклонений и их связь с опорными значениями.
В сравнении 2 доля положительных элементов составляет $0.27$--$0.34$,
в остальных сравнениях она находится в пределах $0.49$--$0.522$
(табл.~\ref{tab:diag}). По принятому порогу выделяется сравнение 2.
Без доли нулевых отклонений $b_0$ это не устанавливает дисбаланс
положительных и отрицательных значений. Корреляции $c$ также
не устанавливают независимости или нулевого
среднего ошибок. Поэтому эти характеристики дополняют норму,
но не заменяют проверку предпосылок~\eqref{eq:moments}.

\begin{table}[!htbp]
    \centering
    \caption{Диагностические характеристики GEMM; номера пар см. рис.~\ref{fig:configs}.
        Последний столбец: выполнен ли порог $|b-0.5|\ge0.1$.}
    \label{tab:diag}
    \begin{tabular}{c c c c}
        \hline
        \textbf{Пара} & \textbf{$b$} & \textbf{$c$} & \textbf{Порог} \\
        \hline
        1 & $0.498$--$0.522$ & $-0.06$--$0.10$ & нет \\
        2 & $0.27$--$0.34$   & $0.08$--$0.16$  & да \\
        3 & $0.49$--$0.51$   & $-0.06$--$0.10$ & нет \\
        4 & $0.49$--$0.51$   & $-0.06$--$0.10$ & нет \\
        \hline
    \end{tabular}
\end{table}

\subsection{Распространение в прямом проходе}

Локальные измерения показывают, насколько расходятся отдельные слои.
Теперь проверим, сохраняется ли этот масштаб на выходе компонента
после введения измеренных возмущений. В таблице~\ref{tab:rfwd}
итоговое относительное расхождение нормировано на максимальное
локальное по формуле~\eqref{eq:ratio}.
Для сравнений 1 и 4 максимумы равны соответственно $1.71$ и $1.51$;
у RoBERTa получены $0.53$ и $0.43$.
Следовательно, итоговое относительное расхождение может как превышать
максимальное локальное, так и быть ниже него. Сопоставление характеризует
зависимость переноса от компонента, но из-за разных опорных норм
не проверяет условие из утверждения~\ref{st:sink}.

\begin{table}[!htbp]
    \centering
    \caption{Коэффициент накопления $R_{\text{fwd}}$ по группам компонентов
        ($^\dagger$~для 2 приведены оценки порядка величины:
        расхождение сопоставимо с $\varepsilon_{\texttt{FP32}}$).}
    \label{tab:rfwd}
    \fitwidth{%
        \begin{tabular}{l l c c c}
            \hline
            \textbf{Компонент / группа}    & \textbf{Нормализация}           &
            \textbf{1} & \textbf{4} &
            \textbf{2$^\dagger$}                                                                              \\
            \hline
            RoBERTa                        & Pre-Norm LN                     & $\mathbf{0.53}$ & $0.43$         & $1.97$        \\
            LeNet                          & нет                             & $1.08$          & $0.81$         & $1.29$        \\
            Механизм внимания (8)               & RMSNorm                         & $1.42$--$1.46$  & $1.09$--$1.30$ & $4.51$--$4.64$ \\
            RMSNorm FFN                    & RMSNorm                         & $1.47$          & $1.32$         & $4.69$        \\
            SwiGLU / GeGLU FFN             & RMSNorm                         & $1.70$--$1.71$  & $1.48$--$1.51$ & $4.05$--$4.07$        \\
            DeepNorm                       & Post-LN                         & $1.39$          & $1.15$         & $4.15$        \\
            \hline
        \end{tabular}%
    }
\end{table}

В сравнении 2 локальное расхождение сопоставимо с машинным эпсилоном
$\varepsilon_{\texttt{FP32}}$, а коэффициенты RMSNorm-компонентов составляют $4.05$--$4.69$.
Эти отношения следует рассматривать как наблюдения, требующие проверки
устойчивости локального знаменателя. Связь с долей $b$ не устанавливает
причины переноса: средние,
направления и совместные ковариации из~\eqref{eq:moments}
здесь не восстановлены.

\subsection{Сопоставление прямого и обратного проходов}

Чтобы сравнить направления переноса, на рис.~\ref{fig:backward-results}
представлены $R_{\mathrm{fwd}}$ и $R_{\mathrm{bwd}}$ для одной пары
реализаций~--- сравнения 1. Обратный проход выполняется на опорных
тензорах прямого прохода, поэтому его ошибки исследуются отдельно.
В прямом проходе коэффициенты находятся в пределах $0.53$--$1.71$,
в обратном~--- $0.35$--$1.87$. Это отношения относительных расхождений
со своими опорными нормами; они не являются нормами якобианов.

\begin{figure}[htbp]
    \centering
    \begin{tikzpicture}[font=\normalsize, x=2.1cm, y=0.7cm]
        \node[align=center] at (1,1.2) {Прямой проход\\$R_{\mathrm{fwd}}$};
        \node[align=center] at (3.65,1.2) {Обратный проход\\$R_{\mathrm{bwd}}$};
        \foreach \offset in {0,2.65} {
            \foreach \tick in {0,0.5,1,1.5,2} {
                \draw[black!15] (\offset+\tick,0.3) -- (\offset+\tick,-5.3);
                \draw (\offset+\tick,-5.6) -- (\offset+\tick,-5.7);
                \node[below] at (\offset+\tick,-5.7) {$\tick$};
            }
            \draw[line width=0.6pt] (\offset,-5.6) -- (\offset+2,-5.6);
        }
        \node[anchor=east] at (-0.12,0) {LeNet};
        \node[draw=blue!65!black, circle, fill=white, inner sep=1.8pt] at (1.08,0) {};
        \node[draw=orange!80!black, circle, fill=white, inner sep=1.8pt] at (4.17,0) {};
        \node[anchor=east] at (-0.12,-1) {RoBERTa};
        \node[draw=blue!65!black, circle, fill=white, inner sep=1.8pt] at (0.53,-1) {};
        \node[draw=orange!80!black, circle, fill=white, inner sep=1.8pt] at (4.41,-1) {};
        \node[anchor=east] at (-0.12,-2) {Механизм внимания};
        \draw[blue!65!black, line width=1pt] (1.42,-2) -- (1.46,-2);
        \draw[blue!65!black] (1.42,-2.12) -- (1.42,-1.88);
        \draw[blue!65!black] (1.46,-2.12) -- (1.46,-1.88);
        \draw[orange!80!black, line width=1pt] (4.34,-2) -- (4.52,-2);
        \draw[orange!80!black] (4.34,-2.12) -- (4.34,-1.88);
        \draw[orange!80!black] (4.52,-2.12) -- (4.52,-1.88);
        \node[anchor=east] at (-0.12,-3) {SwiGLU / GeGLU};
        \draw[blue!65!black, line width=1pt] (1.70,-3) -- (1.71,-3);
        \draw[blue!65!black] (1.70,-3.12) -- (1.70,-2.88);
        \draw[blue!65!black] (1.71,-3.12) -- (1.71,-2.88);
        \draw[orange!80!black, line width=1pt] (4.04,-3) -- (4.06,-3);
        \draw[orange!80!black] (4.04,-3.12) -- (4.04,-2.88);
        \draw[orange!80!black] (4.06,-3.12) -- (4.06,-2.88);
        \node[anchor=east] at (-0.12,-4) {RMSNorm FFN};
        \node[draw=blue!65!black, circle, fill=white, inner sep=1.8pt] at (1.47,-4) {};
        \node[draw=orange!80!black, circle, fill=white, inner sep=1.8pt] at (4.12,-4) {};
        \node[anchor=east] at (-0.12,-5) {DeepNorm};
        \node[draw=blue!65!black, circle, fill=white, inner sep=1.8pt] at (1.39,-5) {};
        \node[draw=orange!80!black, circle, fill=white, inner sep=1.8pt] at (3.00,-5) {};
    \end{tikzpicture}
    \caption{Коэффициенты распространения при раздельном исследовании
        прямого и обратного проходов для сравнения 1.
        Точки соответствуют отдельным компонентам, отрезки~---
        диапазонам компонентов группы, а не доверительным интервалам.}
    \label{fig:backward-results}
\end{figure}
